\section{The Four Population Definitions: Characteristics and Divergence}
\label{sec:four_definitions}

\subsection{Total vs. Total VAP}

The difference between total population and total VAP is the minor population (persons under 18). Minors constitute approximately 22.3\% of the 2020 Census total population (73.1 million out of 331.4 million). The geographic distribution of minors is not uniform: counties and census tracts with younger age profiles (many immigrant communities, rural agricultural areas) have higher minor fractions (25--30\%); counties with older profiles (retirement communities, dense urban cores) have lower minor fractions (15--18\%).

The redistricting consequence of switching from total population to total VAP is modest but directionally consistent: Hispanic-majority border communities, which have high minor fractions, lose relative population weight under total-VAP redistricting. Non-Hispanic white majority communities, which have lower minor fractions (due to older age distribution), gain relative weight. The partisan direction of this shift is toward Republican districts, though the magnitude is smaller than the total-vs.-CVAP shift documented below.

\subsection{Total vs. Citizen VAP}

The difference between total population and citizen VAP reflects both age structure and citizenship composition. Noncitizen adults constitute approximately 10\% of the total national VAP. As documented in N.3~\citep{dellalibera2026n3}, the noncitizen fraction exceeds 15\% in California, Texas, New York, and Florida, with concentrations exceeding 25\% in specific urban tracts.

The divergence between total population and citizen VAP at the congressional district level---the relevant unit for redistricting---averages 12.3 percentage points in the 80 highest-divergence districts, all of which are urban-Democratic districts in the four high-immigration states.

\subsection{Total vs. Adjusted Total}

The difference between total population and adjusted total reflects the group-quarters corrections:
\begin{itemize}
    \item \textbf{Prison adjustment}: Approximately 2.1 million incarcerated persons (2020 Group Quarters types 101--103) are reallocated from facility tracts to estimated home-county tracts. Prison-dense rural counties lose population; urban counties (disproportionately from which prisoners originate) gain population.
    \item \textbf{College adjustment} (optional): Approximately 8.6 million college students in campus group quarters (GQ type 700-series) are reallocated to home-state counties. College-town counties lose population; suburban counties gain population.
\end{itemize}

The combined prison-plus-college adjustment is analyzed separately from the prison-only adjustment in Section~\ref{sec:tract_reassignment}. The primary recommendation in Section~\ref{sec:recommendation} concerns prison adjustment only, consistent with Census Bureau plans for the 2030 redistricting file.

\subsection{Summary of Divergence Magnitude}

\begin{table}[htbp]
\centering
\begin{tabular}{lcccc}
\toprule
\textbf{Comparison} & \textbf{National Mean Reassignment} & \textbf{Max State Reassignment} & \textbf{Direction} \\
\midrule
Total vs. Total VAP & 0.4\% & 1.1\% (TX) & R-leaning \\
Total vs. Citizen VAP & 3.2\% & 11.8\% (CA) & R-leaning \\
Total vs. Adj. Total (prison only) & 0.8\% & 2.1\% (TX) & D-leaning \\
Total vs. Adj. Total (prison+college) & 1.4\% & 3.3\% (IL) & Mixed \\
\bottomrule
\end{tabular}
\caption{Summary of tract reassignment rates and partisan direction across the four population definition comparisons. National mean is population-weighted across 43 redistricting states (excluding single-seat states). Direction indicates which party's districts benefit from the switch, relative to the total-population baseline.}
\label{tab:divergence_summary}
\end{table}

The citizen-VAP comparison produces by far the largest redistricting disruption: 3.2\% mean tract reassignment nationally, reaching nearly 12\% in California. Prison adjustment produces a more modest 0.8\% mean reassignment. The total-vs.-total-VAP comparison produces the smallest disruption at 0.4\%.
